\section*{第 \ref{ch:b1:elimination} 章——\texorpdfstring{$\beta$}{贝塔}-消除}

\begin{solution}{exo:b1:elimination:1}
2-溴丁烷：1-丁烯、\cip{E}-和 \cip{Z}-2-丁烯；主产物
为 \cip{E}-2-丁烯。2-溴-2-甲基丁烷：2-甲基-2-丁烯（主产物，
三取代）和 2-甲基-1-丁烯。溴环己烷：仅得环己烯。
\end{solution}

\begin{solution}{exo:b1:elimination:2}
E1：$v = k[\mathrm{RBr}]$；E2：$v = k[\mathrm{RBr}][\ce{EtO-}]$。测定
不同乙醇钠浓度下生成烯烃的\omterm{def:b1:rate-laws:initial-rate}{初始速率}：
E2 时与之成正比增大，E1 时保持不变。
\end{solution}

\begin{solution}{exo:b1:elimination:3}
沿 C2--C3 看，C2 上的 Br 可与 C3 的任一个 H 反式，也可与
C4 上的甲基反式：两种与 Br 反式的 H \omterm{def:b1:stereochemistry-in-depth:isomers}{构象}可通过 E2
反应（生成 \cip{E} 和 \cip{Z} 烯烃）；第三种\omterm{def:b1:stereochemistry-in-depth:isomers}{构象}不能
向 C3 方向消除。
\end{solution}

\begin{solution}{exo:b1:elimination:4}
SN2（伯卤代烷，氢氧根，低温）；E2（叔卤代烷，强位阻
碱）；SN2（碘负离子，良好\omterm{def:b1:electronic-effects:nucleophile}{亲核试剂}、\omterm{def:b1:predominant-reaction:strong-acid}{弱碱}，\omterm{def:b1:intermolecular-forces-solvents:solvent-classes}{非质子性溶剂}）；
SN1 与 E1，升温时 E1 所占比例增大。
\end{solution}

\begin{solution}{exo:b1:elimination:5}
C3 上只有一个氢。在每个\omterm{def:b1:stereochemistry-in-depth:enantiomers}{非对映异构体}中，使该 H 与 C2 的 Br 反式，
即确定了其余基团的位置：C2 的甲基与 C3 的乙基
在一种\omterm{def:b1:stereochemistry-in-depth:enantiomers}{非对映异构体}中同侧，在另一种中异侧。其中一种
生成 \cip{E}-3-甲基-2-戊烯，另一种生成
\cip{Z}-3-甲基-2-戊烯：\omterm{def:b1:nucleophilic-substitution:stereospecific}{立体专一反应}。
\end{solution}

\begin{solution}{exo:b1:elimination:6}
羟基质子化生成 \ce{R-OH2+}；水离去，得到叔\omterm{def:b1:electronic-effects:intermediates}{碳正离子}
\ce{(CH3)2C+-CH2CH3}；碱（水、\ce{HSO4-}）从相邻碳上夺去
一个质子。主产物（\omterm{prop:b1:elimination:zaitsev}{札依采夫规则}）：
2-甲基-2-丁烯。
\end{solution}

\begin{solution}{exo:b1:elimination:7}
叔丁基使\omterm{def:b1:stereochemistry-in-depth:conformations}{椅式构象}中自身处于平伏位。在
\textit{反式}异构体中 Br 同样处于平伏位，从不与
C--H 反式，E2 只能等待极不利的翻转\omterm{def:b1:stereochemistry-in-depth:conformations}{椅式构象}。
\textit{顺式}异构体中 Br 处于\omterm{def:b1:stereochemistry-in-depth:conformations}{直立键}，且有两个
反式双\omterm{def:b1:stereochemistry-in-depth:conformations}{直立键}氢，E2 很快进行。
\end{solution}

\begin{solution}{exo:b1:elimination:8}
乙醇钠体积小，夺去可生成取代程度更高的
烯烃的氢（\omterm{prop:b1:elimination:zaitsev}{札依采夫规则}）。叔丁醇钠体积大，较易接近甲基上的
暴露氢，而难以接近内部的
\ce{CH2}上的氢，故生成取代程度较低的烯烃。
\end{solution}

\begin{solution}{exo:b1:elimination:9}
溴甲烷没有 $\beta$ 碳。在 1-溴-2,2-二甲基丙烷中，
$\beta$ 碳为季碳，不带氢。
\end{solution}

\begin{solution}{exo:b1:elimination:10}
两种产物都需要\omterm{def:b1:electronic-effects:intermediates}{碳正离子}，它在慢步骤中生成，速率为
$k_1[\mathrm{RBr}]$。随后它在被乙醇捕获（\omterm{def:b1:rate-laws:order}{速率常数} $k_S$）
与失去质子（\omterm{def:b1:rate-laws:order}{速率常数} $k_E$）之间分配，二者
对溶剂均为准一级，烯烃的比例为
$k_E/(k_E + k_S)$，与\omterm{def:b1:nucleophilic-substitution:substitution}{底物}浓度无关。
\end{solution}

\begin{solution}{exo:b1:elimination:11}
从两个 Br 呈反式的\omterm{def:b1:stereochemistry-in-depth:enantiomers}{内消旋化合物}\omterm{def:b1:stereochemistry-in-depth:isomers}{构象}出发
（此时两个苯基反式，两个 H 反式），将 C1 旋转 $120^\circ$，
使其 H 与 C2 的 Br 反式。此时两个苯基位于
H--Br 轴的同侧，二者最终为\textit{顺式}。C1 上 Br 优先于
Ph；C2 上 Ph 优先于 H；C2 的 Br 与苯基为\textit{反式}：
\cip{E}-1-溴-1,2-二苯基乙烯。
\end{solution}

\begin{solution}{exo:b1:elimination:12}
只有\omterm{def:b1:electronic-effects:nucleophile}{离去基团}处于\omterm{def:b1:stereochemistry-in-depth:conformations}{直立键}的\omterm{def:b1:stereochemistry-in-depth:conformations}{椅式构象}才反应：$v = k\,x\,[\mathrm{RY}][\mathrm{B}]$，
其中 $x$ 为该\omterm{def:b1:stereochemistry-in-depth:conformations}{椅式构象}所占分数。它所需的每个直立烷基
需付出若干 kJ/mol 的能量代价（甲基为 \qty{5.7}{kJ/mol}），
室温下每个基团使 $x$ 约减为原来的十分之一。两三个此类基团时，
$x$ 降至 $10^{-2}$ 或 $10^{-3}$，反应也相应慢得多。
\end{solution}

\begin{solution}{pb:b1:elimination:1}
\textbf{1.} Cl、异丙基和甲基均处于\omterm{def:b1:stereochemistry-in-depth:conformations}{平伏键}的\omterm{def:b1:stereochemistry-in-depth:conformations}{椅式构象}。
\textbf{2.} 异丙基和甲基处于\omterm{def:b1:stereochemistry-in-depth:conformations}{平伏键}，Cl 处于\omterm{def:b1:stereochemistry-in-depth:conformations}{直立键}。
\textbf{3.} \omterm{def:b1:stereochemistry-in-depth:enantiomers}{非对映异构体}：二者仅在 C1 处不同。
\textbf{4.} \omterm{def:b1:electronic-effects:nucleophile}{离去基团}与 $\beta$ 氢呈反式双\omterm{def:b1:stereochemistry-in-depth:conformations}{直立键}。
\textbf{5.} 只有在翻转后的\omterm{def:b1:stereochemistry-in-depth:conformations}{椅式构象}中，Cl、异丙基和甲基均
处于\omterm{def:b1:stereochemistry-in-depth:conformations}{直立键}。
\textbf{6.} 三个基团均处于\omterm{def:b1:stereochemistry-in-depth:conformations}{直立键}：仅甲基就需付出约
\qty{5.7}{kJ/mol}，较大的异丙基付出更多，氯则较少，
合计约 \qty{13}{kJ/mol}，因此分数约为
$\exp(-13000/2479) = \num{5e-3}$，即题目所取的数值。
\textbf{7.} C2（一个 H，连有异丙基）和 C6（两个 H）。
\textbf{8.} C1 上 Cl 处于\omterm{def:b1:stereochemistry-in-depth:conformations}{直立键}时，C2 的直立 H（异丙基处于
平伏位）和 C6 的直立 H：共两个。
\textbf{9.} 在三直立的\omterm{def:b1:stereochemistry-in-depth:conformations}{椅式构象}中，异丙基占据 C2 的\omterm{def:b1:stereochemistry-in-depth:conformations}{直立键}，
因此该 H 为平伏，不反式；C6 有直立 H，为反式。
\textbf{10.} 新薄荷基氯有两个，薄荷基氯有一个。
\textbf{11.} 沿 C1--C6：C1 上的 Cl（直立，向上）与 C6 上的直立 H（向下）
相对；环键和平伏 H 交叉排列于两者之间。
\textbf{12.} 在\omterm{def:b1:stereochemistry-in-depth:conformations}{椅式构象}中，没有任何 C--H 键与 C--Cl 键\omterm{def:b1:elimination:periplanar}{顺式共平面}；
强行如此则需要重叠\omterm{def:b1:stereochemistry-in-depth:isomers}{构象}，环张力增大。
\textbf{13.} 消去 C2 上的 H：4-甲基-1-（丙-2-基）环己-1-烯
（三取代）；消去 C6 上的 H：3-甲基-6-（丙-2-基）环己-1-烯
（二取代）。
\textbf{14.} 三取代烯烃，取代程度更高且更稳定。
\textbf{15.} 两种烯烃中三取代者占 \qty{78}{\percent}：与\omterm{prop:b1:elimination:zaitsev}{札依采夫规则}一致。
\textbf{16.} 只生成二取代烯烃，与\omterm{prop:b1:elimination:zaitsev}{札依采夫规则}的
预测相反，由唯一可用的反式氢所决定。
\textbf{17.} 两者均为\omterm{def:b1:nucleophilic-substitution:stereospecific}{立体专一反应}（需反式）；新薄荷基氯
发生\omterm{def:b1:elimination:selectivity}{区域选择性反应}（78/22）；薄荷基氯只得
一种\omterm{def:b1:stereochemistry-in-depth:isomers}{构造异构体}。
\textbf{18.} 连 Cl 的碳为\omterm{def:b1:electronic-effects:carbon-class}{仲碳}，旁边紧挨异丙基，
空间位阻大，不利于 SN2；而乙醇钠为\omterm{def:b1:predominant-reaction:strong-acid}{强碱}。
\textbf{19.} $v \propto x\,n_{\mathrm{H}}$：反应性\omterm{def:b1:stereochemistry-in-depth:conformations}{椅式构象}的分数乘以
反式氢的个数。
\textbf{20.} $0.95 \times 2 = 1.90$.
\textbf{21.} $\num{5.0e-3} \times 1 = \num{5.0e-3}$.
\textbf{22.} \textbf{$k(\text{新薄荷基氯})/k(\text{薄荷基氯}) = 1.90/\num{5.0e-3}
= 380$}.
\end{solution}
